% Auto-generated from report/3.md — do not edit by hand.
\section{مجموعه‌های داده معیار و تحلیل ساختاری آن‌ها}\label{sec:3-1}


برای ارزیابی مدل‌های پیشنهادی و مقایسه با خطوط مبنا، از چهار پایگاه داده استاندارد زیست‌مولکولی استفاده شده است. این داده‌ها طیف کاملی از ساختارهای شبکه‌ای را پوشش می‌دهند: گراف‌های همگن و دوقطبی، گراف‌های متراکم و تنک، و مقیاس‌های مختلف از ۱,۵۱۴ تا ۱۹,۷۸۳ گره. تمامی داده‌ها از پوشه \pcode{fold1} مخزن رسمی \lr{SkipGNN} \cite{huang2020skipgnn} استخراج شده‌اند و بدون هیچ تغییری در تقسیم‌بندی استفاده می‌شوند.


\begin{table}[htbp]
\centering
\footnotesize
\caption{مشخصات آماری و توپولوژیک چهار بنچمارک استاندارد}
\label{tab:3-1}
\begin{adjustbox}{max width=\textwidth}
\begin{tabular}{lllllllll}
\toprule
مجموعه داده & نوع شبکه & نود مبدا & نود مقصد & کل نودها ($N$) & یال آموزش & یال اعتبارسنجی & یال آزمون & میانگین درجه \\
\midrule
\lr{DDI} & همگن & ۱,۵۱۴ (دارو) &  -  & ۱,۵۱۴ & ۶۷,۹۱۹ & ۹,۷۰۳ & ۱۹,۴۰۶ & ۶۴.۰۹ \\
\lr{PPI} & همگن & ۵,۶۰۴ (پروتئین) &  -  & ۵,۶۰۴ & ۳۲,۶۵۱ & ۴,۶۶۴ & ۹,۳۲۹ & ۸.۳۲ \\
\lr{DTI} & دوقطبی & ۵,۰۱۸ (دارو) & ۲,۳۲۵ (پروتئین) & ۷,۳۴۳ & ۲۱,۱۹۴ & ۳,۰۲۸ & ۶,۰۵۶ & ۴.۱۲ \\
\lr{GDI} & دوقطبی & ۹,۴۱۳ (ژن) & ۱۰,۳۷۰ (بیماری) & ۱۹,۷۸۳ & ۱۱۴,۴۴۵ & ۱۶,۳۴۹ & ۳۲,۶۹۸ & ۸.۲۶ \\
\bottomrule
\end{tabular}
\end{adjustbox}
\end{table}


\subsection{پایگاه داده برهم‌کنش دارو-دارو (\lr{BioSNAP-DDI})}\label{sec:3-1-1}


این شبکه شامل تداخلات فارماکوکینتیکی و فارماکودینامیکی بین داروهای تأییدشده توسط \lr{FDA} است که از پایگاه داده \lr{DrugBank} استخراج و توسط مجموعه \lr{BioSNAP} استنفورد توزیع شده است \cite{bionsnap}.

\textbf{ویژگی‌های ساختاری:}
\begin{itemize}
\item گراف همگن متقارن با $N = 1{,}514$ دارو و $67{,}919$ یال آموزشی مثبت.
\item میانگین درجه بسیار بالا ($\bar{d} = 64.09$): متراکم‌ترین گراف در بنچمارک.
\item ساختار جوامع هم‌پوشان غالب است. در چنین ساختاری، مسیرهای ۲-گام هم‌نوع دارای غنای اطلاعاتی فراوانی هستند.
\end{itemize}

\subsection{پایگاه داده برهم‌کنش پروتئین-پروتئین (\lr{HuRI-PPI})}\label{sec:3-1-2}


اینتراکتوم مرجع پروتئینی انسان \cite{luck2020huri} که از طریق آزمایش‌های غربال‌گری دو رگه‌ای مخمری با توان بالا به دست آمده است.

\textbf{ویژگی‌های ساختاری:}
\begin{itemize}
\item گراف همگن متقارن با $N = 5{,}604$ پروتئین و $32{,}651$ یال آموزشی مثبت.
\item میانگین درجه تنک ($\bar{d} = 8.32$).
\item ساختار مدولار زیستی با ضریب خوشه‌بندی محلی بالا، بازتاب‌دهنده کمپلکس‌های پروتئینی سلولی.
\end{itemize}

\subsection{پایگاه داده برهم‌کنش دارو-هدف (\lr{BioSNAP-DTI})}\label{sec:3-1-3}


تعاملات مهاری و اتصالی مولکول‌های دارو به پروتئین‌های هدف انسانی، استخراج‌شده از داده‌های \lr{ChG-Miner} و \lr{DrugBank}.

\textbf{ویژگی‌های ساختاری:}
\begin{itemize}
\item گراف دوقطبی با $N = 7{,}343$ موجودیت: $N_{\text{\lr{drug}}} = 5{,}018$ دارو ($\mathcal{V}_1$) و $N_{\text{\lr{protein}}} = 2{,}325$ پروتئین ($\mathcal{V}_2$).
\item $21{,}194$ یال آموزشی مثبت با میانگین درجه بسیار تنک ($\bar{d} = 4.12$).
\item نسبت داروها به پروتئین‌ها حدود $2.15$ است.
\end{itemize}

\subsection{پایگاه داده برهم‌کنش ژن-بیماری (\lr{DisGeNET-GDI})}\label{sec:3-1-4}


پایگاه داده جامع \lr{DisGeNET} \cite{pinero2020disgenet} که همبستگی‌های تجربی میان ژن‌های انسانی و فنوتیپ‌های بیماری را بر پایه شواهد بالینی و ژنتیکی گردآوری می‌کند.

\textbf{ویژگی‌های ساختاری:}
\begin{itemize}
\item گراف دوقطبی کلان با $N = 19{,}783$ موجودیت: $N_{\text{\lr{gene}}} = 9{,}413$ ژن و $N_{\text{\lr{disease}}} = 10{,}370$ بیماری.
\item $114{,}445$ یال آموزشی مثبت با میانگین درجه $8.26$.
\item بزرگ‌ترین چالش محاسباتی و حافظه‌ای را به دلیل ابعاد ماتریس مجاورت ایجاد می‌کند.
\end{itemize}

\section{پایپ‌لاین بدون نشت داده (\lr{Leak-Free Pipeline})}\label{sec:3-2}


یکی از مهم‌ترین نیازمندی‌های روش‌شناختی در پیش‌بینی پیوند، تضمین عدم نشت اطلاعات از مجموعه‌های اعتبارسنجی و آزمون به ساختار گراف آموزشی است. پایپ‌لاین پیاده‌سازی‌شده در تابع \pcode{load\_dataset\_splits} (ماژول \pcode{src/data/loader.py}) این تضمین را در پنج مرحله فراهم می‌کند:

\textbf{مرحله ۱  -  نگاشت شناسه‌ها:} رشته‌های شناسایی موجودیت‌ها به اندیس‌های صحیح پیوسته نگاشت می‌شوند. در گراف‌های دوقطبی، تابع \pcode{\_build\_idx\_map} با استفاده از هیوریستیک نوع‌شناسی (پیشوند \pcode{DB} برای داروها در \lr{DTI}، شناسه عددی برای ژن‌ها در \lr{GDI})، منابع و اهداف را به‌طور قطعی تفکیک می‌کند. در صورت انحراف بیش از یک واحد از تعداد مورد انتظار منابع، لودر خطا صادر می‌کند.

\textbf{مرحله ۲  -  جهت‌یابی دوقطبی:} در گراف‌های \lr{DTI} و \lr{GDI}، تابع \pcode{\_orient\_bipartite} تضمین می‌کند که ستون اول همواره مبدا و ستون دوم همواره مقصد باشد:


\begin{equation}
\label{eq:3.1}
\text{\lr{if }} u \ge N_{\text{\lr{src}}} \land v < N_{\text{\lr{src}}} \implies \text{\lr{swap}}(u, v)
\end{equation}


\textbf{مرحله ۳  -  ساخت ماتریس مجاورت آموزشی:} ماتریس $A_{\text{\lr{train}}}$ صرفاً از یال‌های مثبت آموزشی ساخته می‌شود:


\begin{equation}
\label{eq:3.2}
\mathcal{E}_{\text{\lr{train}}} = \{(u, v) \in \mathcal{D}_{\text{\lr{train}}} : y_{uv} = 1\}
\end{equation}



\begin{equation}
\label{eq:3.3}
A_{\text{\lr{train}}} = \text{\lr{SymmetricClosure}}(\mathcal{E}_{\text{\lr{train}}})
\end{equation}


\textbf{مرحله ۴  -  ساخت عملگرهای پرش:} تمام ماتریس‌های پرش ($F_s$, $W_2$, $W_3$) مستقیماً از $A_{\text{\lr{train}}}$ ساخته می‌شوند. هیچ ردی از یال‌های اعتبارسنجی یا آزمون در آن‌ها وجود ندارد.

\textbf{مرحله ۵  -  ثبت مثبت‌های شناخته‌شده:} یال‌های مثبت اعتبارسنجی و آزمون در یک مجموعه هش (\pcode{known\_positives}) نگهداری می‌شوند تا در مرحله نمونه‌برداری منفی، از انتخاب تصادفی یال‌های واقعی جلوگیری شود. این یال‌ها \textbf{هرگز} وارد $A_{\text{\lr{train}}}$ نمی‌شوند. صحت این تضمین در آزمون‌های واحد \pcode{test\_ddi\_shapes\_and\_no\_leakage} و \pcode{test\_dti\_bipartite\_ranges} اعتبارسنجی شده است.

\subsection{انتخاب آگاهانه ویژگی‌های ورودی}\label{sec:3-2-1}


در این پژوهش، ویژگی‌های ورودی گره‌ها به ماتریس هویت تنک ($X = I_N$) محدود شده‌اند. این یک تصمیم روش‌شناختی آگاهانه است:

\textbf{هدف:} سنجش ظرفیت خالص معماری‌های پیام‌رسانی در کشف الگوهای توپولوژیک پیوند، بدون دخالت ویژگی‌های بیوشیمیایی.

\textbf{توجیه:} افزودن ویژگی‌های غنی (مانند اثرانگشت مورگان یا امبدینگ‌های \lr{ESM-2}) می‌تواند اثر معماری را با اثر ویژگی‌ها مخلوط کند و تفکیک سهم هر یک را دشوار سازد. این محدودیت به‌صراحت در فصل ششم به‌عنوان مسیر کار آینده ذکر شده است.

\subsection{ساخت تنسورهای تنک کوالس‌شده}\label{sec:3-2-2}


تمام ماتریس‌های مجاورت ابتدا به فرمت تنک فشرده سطری (\pcode{scipy.sparse.csr\_matrix}) ساخته شده و پس از اعمال عملگرهای ریاضی، به فرمت مختصات تنک ۳۲ بیتی تبدیل می‌شوند. تابع \pcode{scipy\_to\_torch\_sparse} تنسور تنک ادغام‌شده (\pcode{torch.sparse\_coo\_tensor().coalesce()}) را مستقیماً روی حافظه \lr{GPU} بارگذاری می‌کند.

\section{عملگرهای پرش توپولوژیک (\lr{Topological Skip Operators})}\label{sec:3-3}


سه دسته عملگر پرش توپولوژیک به صورت ماتریس‌های تنک بهینه‌شده طراحی و پیاده‌سازی شده‌اند (ماژول \pcode{src/data/normalization.py}):

\subsection{عملگر پرش باینری بدون سلف‌لوپ (\lr{Binary Skip})}\label{sec:3-3-1}


بازتولید ساختار پرش مقاله مرجع \cite{huang2020skipgnn} با یک اصلاح ساختاری: قطر اصلی صراحتاً صفر می‌شود:


\begin{equation}
\label{eq:3.4}
A_s^{\text{\lr{bin}}} = \text{\lr{sign}}(A_{\text{\lr{train}}} \cdot A_{\text{\lr{train}}}^T), \quad (A_s^{\text{\lr{bin}}})_{uu} = 0 \quad \forall u
\end{equation}


\subsection{عملگر پرش پیوسته تخصیص منابع (\lr{Resource Allocation Skip})}\label{sec:3-3-2}


برای حفظ اطلاعات شدت پیوند و چگالی همسایگان مشترک، عملگر تخصیص منابع \cite{zhou2009predicting} جایگزین عملگر علامت می‌شود:


\begin{equation}
\label{eq:3.5}
(D^{-1})_{kk} = \begin{cases} \frac{1}{d_k} & \text{\lr{if }} d_k > 0 \\ 0 & \text{\lr{if }} d_k = 0 \end{cases}
\end{equation}



\begin{equation}
\label{eq:3.6}
W_2 = A_{\text{\lr{train}}} \cdot D^{-1} \cdot A_{\text{\lr{train}}}^T, \quad (W_2)_{uu} = 0 \quad \forall u
\end{equation}


عنصر $(W_2)_{uv}$ بیانگر مجموع وزن‌های اشتراک میان دو گره $u$ و $v$ است:


\begin{equation}
\label{eq:3.7}
(W_2)_{uv} = \sum_{k \in \mathcal{N}(u) \cap \mathcal{N}(v)} \frac{1}{d_k}
\end{equation}


\textbf{تحلیل رفتار:} اگر دو دارو از طریق یک پروتئین اختصاصی با درجه پایین ($d_k = 2$) متصل باشند، وزن پرش $0.5$ خواهد بود؛ در حالی که اگر از طریق یک پروتئین هاب غیراختصاصی ($d_k = 100$) متصل باشند، وزن پرش به $0.01$ کاهش می‌یابد. این فرمولاسیون به طور خودکار نویز هاب‌ها را مهار می‌کند.

\textbf{نکته پیاده‌سازی:} در کد (\pcode{build\_weighted\_skip})، برای جلوگیری از تقسیم بر صفر، از $\max(1, d_k)$ استفاده شده است. این معادل همان رفتار ریاضی معادله (۳-۵) است، زیرا گره‌های با درجه صفر در $A_{\text{\lr{train}}}$ هیچ همسایه‌ای ندارند و در مجموع ظاهر نمی‌شوند.

\subsection{عملگر مسیر بازگشت ۳-گام (\lr{Bipartite 3-Walk Return})}\label{sec:3-3-3}


همان‌طور که در بخش ۲-۴-۱ اثبات شد، در گراف‌های دوقطبی، ماتریس $W_2$ اتصالات درون‌نوعی را برقرار می‌کند. برای انتقال پیام به موجودیت هدف کاندید، یک گام انتشار اضافی لازم است:


\begin{equation}
\label{eq:3.8}
W_3 = W_2 \cdot D^{-1} \cdot A_{\text{\lr{train}}}
\end{equation}


\textbf{مفهوم زیستی:} در گراف \lr{DTI}، مقدار $(W_3)_{dp}$ نشان‌دهنده امتیاز انتقال پیام از داروی $d$ به پروتئین $p$ از طریق تمام مسیرهای طول ۳ ($d \to p' \to d' \to p$) است.

\subsection{نرمال‌سازی لاپلاسی متقارن ایمن (\lr{Safe Laplacian Normalization})}\label{sec:3-3-4}


در پیاده‌سازی‌های رایج، برای جلوگیری از تقسیم بر صفر، مقدار کوچک $\epsilon$ اضافه می‌شود. اشکال ریاضی: اگر نودی ایزوله باشد ($d = 0$)، اضافه کردن $\epsilon = 10^{-5}$ باعث می‌شود مقدار نرمال‌شده $\frac{1}{\sqrt{10^{-5}}} \approx 316$ شود.

\textbf{راهکار پیاده‌سازی‌شده} (تابع \pcode{laplacian\_normalize}):


\begin{equation}
\label{eq:3.9}
(D^{-1/2})_{ii} = \begin{cases} d_i^{-0.5} & \text{\lr{if }} d_i > 0 \\ 0 & \text{\lr{otherwise}} \end{cases}
\end{equation}



\begin{equation}
\label{eq:3.10}
F = D^{-1/2} A D^{-1/2}
\end{equation}


برای گراف اصلی، ابتدا سلف‌لوپ اضافه می‌شود ($A + I_N$). برای گراف‌های پرش، بدون سلف‌لوپ محاسبه می‌گردد. این رفتار در آزمون واحد \pcode{test\_laplacian\_isolated\_nodes\_are\_zero} اعتبارسنجی شده است.


\begin{table}[htbp]
\centering
\footnotesize
\caption{خلاصه عملگرهای توپولوژیک و مدل‌های مصرف‌کننده}
\label{tab:3-2}
\begin{tabular}{llll}
\toprule
تنسور در کد & عملگر پایه & فرمولاسیون لاپلاسی & مدل‌های مصرف‌کننده \\
\midrule
\pcode{f\_orig} & $A + I_N$ & $\tilde{D}^{-1/2}(A+I)\tilde{D}^{-1/2}$ & تمام مدل‌های عصبی \\
\pcode{f\_skip\_bin} & $\text{\lr{sign}}(AA^T) - \text{\lr{diag}}$ & $\tilde{D}_s^{-1/2} A_s^{\text{\lr{bin}}} \tilde{D}_s^{-1/2}$ & \lr{SkipGNNBaseline} \\
\pcode{f\_skip\_weighted} & $AD^{-1}A^T - \text{\lr{diag}}$ & $\tilde{D}_s^{-1/2} W_2 \tilde{D}_s^{-1/2}$ & \lr{AMS}, \lr{GAT}, \lr{Contrastive}, \lr{3-hop} \\
\pcode{f\_3hop} & $W_2 D^{-1} A$ & $\tilde{D}_3^{-1/2} W_3 \tilde{D}_3^{-1/2}$ & \lr{ThreeHopSkipGNN} \\
\bottomrule
\end{tabular}
\end{table}


\section{پروتکل ارزیابی دوگانه (\lr{Dual-Bank Evaluation Protocol})}\label{sec:3-4}


برای سنجش عادلانه و فاقد سوگیری کارایی مدل‌ها، ارزیابی تحت پروتکل دوگانه صورت می‌پذیرد. این پروتکل رفتار مدل را هم در شرایط استاندارد و هم در سخت‌ترین سناریوهای آزمون تنش ارزیابی می‌کند.

\subsection{بانک آزمون یکنواخت (\lr{Uniform Test Bank})}\label{sec:3-4-1}


در این بانک، مجموعه آزمون رسمی (\pcode{test.csv}) بدون تغییر استفاده می‌شود. هدف: ایجاد شرایط کاملاً یکسان با آزمون‌های مقاله مرجع \cite{huang2020skipgnn} به منظور مقایسه مستقیم.

\subsection{بانک آزمون سخت آگاه از نوع موجودیت (\lr{Hard Test Bank})}\label{sec:3-4-2}


هدف: حذف امکان بهره‌برداری مدل از محبوبیت درجه گره‌ها و اجبار شبکه به یادگیری ساختار واقعی پیوند. این رویکرد با توصیه‌های لی و همکاران \cite{li2023evaluating} درباره محدودیت‌های منفی‌های تصادفی هم‌راستا است.

\textbf{الگوریتم} (تابع \pcode{generate\_bipartite\_aware\_hard\_negatives} در \pcode{src/data/samplers.py}):

\textbf{گام ۱  -  محاسبه درجات:} درجات گره‌های تارگت در گراف آموزش محاسبه می‌شود:


\begin{equation}
\label{eq:3.11}
d_v = \sum_u (A_{\text{\lr{train}}})_{uv}
\end{equation}


\textbf{گام ۲  -  افراز چارک‌های درجه‌ای:} گره‌های تارگت بر اساس صدک‌های ۲۵، ۵۰ و ۷۵ درجات به ۴ باکت تقسیم می‌شوند:


\begin{equation}
\label{eq:3.12}
\mathcal{B}_0 = \{v : d_v \le q_{25}\}, \quad \mathcal{B}_1 = \{v : q_{25} < d_v \le q_{50}\}
\end{equation}



\begin{equation}
\label{eq:3.13}
\mathcal{B}_2 = \{v : q_{50} < d_v \le q_{75}\}, \quad \mathcal{B}_3 = \{v : d_v > q_{75}\}
\end{equation}


\textbf{گام ۳  -  نمونه‌برداری هم‌چارک:} به ازای هر یال مثبت آزمون $(u, v)$:
\begin{itemize}
\item چارک درجه‌ای تارگت واقعی $v$ تعیین می‌شود: $b = \text{\lr{Bucket}}(v)$.
\item کاندید منفی $v^-$ از همان چارک انتخاب می‌شود: $v^- \in \mathcal{B}_b$.
\item شروط اعتبار بررسی می‌شوند:
\item در گراف دوقطبی: $(u, v^-) \notin \mathcal{E}_{\text{\lr{known}}}$ و $u \ne v^-$ و $v^-$ از نوع تارگت باشد.
\item در گراف همگن: $(u, v^-) \notin \mathcal{E}_{\text{\lr{known}}}$ و $(v^-, u) \notin \mathcal{E}_{\text{\lr{known}}}$ و $u \ne v^-$.
\end{itemize}

\textbf{گام ۴  -  مکانیزم فال‌بک:} در صورت عدم یافتن کاندید معتبر پس از ۱۰۰ تلاش در همان چارک، فال‌بک با سقف ۱,۰۰۰ تلاش در فضای کل تارگت‌ها اعمال می‌شود. این مکانیزم برای جلوگیری از حلقه‌های بی‌پایان در نودهایی با درجه اشباع طراحی شده است. صحت رفتار این الگوریتم در آزمون واحد \pcode{test\_hard\_negatives\_respect\_bipartite\_types} اعتبارسنجی شده است.

\subsection{اصول حاکم بر پروتکل سخت}\label{sec:3-4-3}


\textbf{پایستگی برش دوقطبی (\lr{Bipartite Cut Invariance}):} در گراف‌های \lr{DTI} و \lr{GDI}، دارو هرگز با داروی دیگر جفت نمی‌شود. نوع موجودیت‌ها حفظ می‌گردد.

\textbf{بررسی تقارن در گراف‌های همگن:} در \lr{DDI} و \lr{PPI}، چون گراف بدون جهت است، باید اطمینان حاصل شود که نه $(u, v^-)$ و نه $(v^-, u)$ در مجموعه تعاملات واقعی قرار ندارند.

\textbf{عدم نشت آستانه:} آستانه تصمیم‌گیری $\tau^*$ صرفاً روی داده‌های اعتبارسنجی محاسبه و سپس منجمد می‌شود (بخش ۵-۱-۳).

\section{معیارهای ارزیابی و انتخاب آستانه}\label{sec:3-5}


ارزیابی نهایی توسط تابع \pcode{evaluate\_dual\_bank} (ماژول \pcode{src/eval/metrics.py}) انجام می‌شود:

\textbf{انتخاب آستانه بهینه:}


\begin{equation}
\label{eq:3.14}
\tau^* = \arg\max_{\tau \in [0.05, 0.95]} \text{\lr{F1}}\left(y_{\text{\lr{val}}}, \mathbb{I}(p_{\text{\lr{val}}} \ge \tau)\right)
\end{equation}


پویش ۹۱ نقطه‌ای در بازه $[0.05, 0.95]$ صرفاً روی احتمالات مجموعه اعتبارسنجی انجام می‌شود. سپس $\tau^*$ برای هر دو بانک آزمون منجمد می‌گردد. صحت این جداسازی در آزمون واحد \pcode{test\_threshold\_is\_validation\_only} تضمین شده است.

\textbf{معیارهای گزارش‌شده:}
\begin{itemize}
\item \textbf{\lr{AUROC}:} مساحت زیر منحنی مشخصه عملکرد گیرنده.
\item \textbf{\lr{AUPRC}:} مساحت زیر منحنی دقت-بازخوانی  -  معیار اصلی رتبه‌بندی در این گراف‌های نامتوازن.
\item \textbf{\lr{F1}:} شاخص تعادلی دقت و بازخوانی در آستانه $\tau^*$.
\item \textbf{\lr{Brier Score}:} میانگین خطای مربعی احتمالات.
\end{itemize}
